\section{Implementation in the workbook}
\label{sec:implementation}

\subsection{Layout}
The calculation lives on the \textbf{Load Flow} tab of the HVDC calculation
workbook. Column~A carries the quantity, column~B its unit, and the three
result columns are the three cases of Table~\ref{tab:cases}: column~C the
symmetrical monopole, column~D the balanced bipole, column~E the bipole with
one pole out. Column~E takes its plant data from column~D by reference, so the
outage case cannot drift out of step with the bipole it belongs to; only the
flow rows differ. Input cells are shaded, results are shaded in the second
house colour, matching the convention already used on the Power and Losses
tabs.

\subsection{Inputs}
\begin{table}[htbp]
  \centering
  \small
  \renewcommand{\arraystretch}{1.25}
  \begin{tabularx}{\linewidth}{@{}c l l X@{}}
    \toprule
    \textbf{Row} & \textbf{Input} & \textbf{Symbol} & \textbf{Remark} \\
    \midrule
    6  & AC power at receiving end   & $P_{ac,inv}$ & At the inverter AC terminals; the design target. \\
    7  & DC voltage pole-to-ground   & $V_{d}$      & At the rectifier DC terminals. \\
    8  & DC line length              & $L$          & Route length, miles. \\
    9--11 & HV conductor, bundles, specific resistance
                                     & $n_{HV}$, $r'_{HV}$
                                     & Per subconductor, $\Omega$/1000\,ft. \\
    12--14 & DMR conductor, bundles, specific resistance
                                     & $n_{DMR}$, $r'_{DMR}$
                                     & Separate input: the DMR is generally not the pole conductor. \\
    15 & Losses inverter             & $p_{inv}$    & Fraction of station throughput. \\
    16 & Losses rectifier            & $p_{rec}$    & Fraction of station throughput. \\
    17 & Margin for DC current       & $m$          & Applied to the rated current only. \\
    \bottomrule
  \end{tabularx}
  \caption{Input rows of the Load Flow tab.}
  \label{tab:inputs}
\end{table}

The DC voltage is an input the original request did not list, but the flow
cannot be determined without it: the same power at the same length gives a
different current, loss and voltage drop at every voltage level, and it is the
voltage that fixes which of them applies. It is entered pole-to-ground, the
same convention as the Losses tab.

\subsection{Formulas as entered}
Table~\ref{tab:formulas} maps every calculated row to the equation it
implements. Cell references are for column~C (symmetrical monopole); columns~D
and~E are the same relations on their own column, with the two differences
noted at the foot of the table.

\begin{table}[htbp]
  \centering
  \small
  \renewcommand{\arraystretch}{1.3}
  \begin{tabularx}{\linewidth}{@{}c X l c@{}}
    \toprule
    \textbf{Row} & \textbf{Formula as entered} & \textbf{Quantity} &
    \textbf{Eq.} \\
    \midrule
    20 & \texttt{=C11*5.28/C10*C8}
       & $R_{HV}$        & \eqref{eq:resistance} \\
    21 & \texttt{=D14*5.28/D13*D8}
       & $R_{DMR}$       & \eqref{eq:resistance} \\
    22 & \texttt{=2*C20}
       & $R_{loop}$      & \eqref{eq:loop} \\
    25 & \texttt{=C6/(1-C15)}
       & $P_{dc,inv}$    & \eqref{eq:step1} \\
    26 & \texttt{=(C7-SQRT(C7\textasciicircum{}2-C22*C25))/C22*1000}
       & $I_{d}$         & \eqref{eq:current} \\
    27 & \texttt{=C26*(1+C17)}
       & $I_{rated}$     & \eqref{eq:margin} \\
    28 & \texttt{=2*C7*C26/1000}
       & $P_{dc,rec}$    & \eqref{eq:c1} \\
    29 & \texttt{=C22*(C26/1000)\textasciicircum{}2}
       & $P_{loss,HV}$   & \eqref{eq:losssplit} \\
    30 & \texttt{0}
       & $P_{loss,DMR}$  & \eqref{eq:losssplit} \\
    31 & \texttt{=C29+C30}
       & $P_{loss,line}$ & \eqref{eq:c2} \\
    32 & \texttt{=C7-C26/1000*C20}
       & $V_{d,inv}$     & \eqref{eq:vinv} \\
    33 & \texttt{=C25*(1-C15)}
       & $P_{ac,inv}$    & \eqref{eq:inv} \\
    34 & \texttt{=C28/(1-C16)}
       & $P_{ac,rec}$    & \eqref{eq:c4} \\
    35 & \texttt{=C34-C33}
       & $P_{loss,total}$& \eqref{eq:c5} \\
    36 & \texttt{=C33/C34*100}
       & $\eta$          & \eqref{eq:c6} \\
    39 & \texttt{=C7\textasciicircum{}2/C22}
       & $P_{max}$       & \eqref{eq:pmax} \\
    \bottomrule
  \end{tabularx}
  \caption{Calculated rows of the Load Flow tab and the equations they
           implement. The $1000$ factors convert between the kiloamperes of
           the equations and the amperes reported in row~26.}
  \label{tab:formulas}
\end{table}

Two rows differ in the outage column, and they are the whole of the difference
between a balanced and an unbalanced bipole:

\begin{itemize}\setlength{\itemsep}{0.3em}
  \item \textbf{Row 22}, loop resistance, becomes \texttt{=E20+E21} --- one pole
        conductor plus the DMR, Equation~\eqref{eq:loop}.
  \item \textbf{Row 26}, current, becomes \texttt{=D26} --- the case is defined
        by holding the pole current, so rows 25 and 28 reverse direction:
        row~28 is \texttt{=E7*E26/1000} (single pole, $k=1$) and row~25 becomes
        \texttt{=E28-E31}, the forward chain of
        Equations~\eqref{eq:m1}--\eqref{eq:m4}.
\end{itemize}

\subsection{Row 33 as a closure check}
In columns C and D, row~33 recomputes $P_{ac,inv}$ from the DC power that
row~25 derived from it. It therefore has to reproduce the input in row~6
exactly:
\[
  \frac{P_{ac,inv}}{1-p_{inv}}\bigl(1-p_{inv}\bigr) = P_{ac,inv}.
\]
Any discrepancy between rows 6 and 33 means a formula in the chain has been
disturbed --- it is a cheap, always-visible check on the sheet. In column~E,
row~33 is not a check but the answer: it is the power the link can still
deliver on one pole.
